%% ----------------------------------------------------------------------------
%%                              Chinese Abstract
%% ----------------------------------------------------------------------------
\begin{abstract}{大语言模型，幻觉检测，内部表征，检索增强生成，稀疏自编码器}

近年来，大语言模型在问答、生成与检索等任务中表现出较强能力，但仍容易产生与事实不符、逻辑错误或缺乏依据的幻觉内容，影响系统的可靠性与安全性。检索增强生成（Retrieval-Augmented Generation, RAG）框架通过引入外部检索知识源，增强了模型的生成能力，一定程度上缓解了幻觉问题；但其并未能根除幻觉，模型仍有可能输出与检索知识不一致的内容。幻觉检测任务旨在有效识别和检测大语言模型生成的幻觉内容，其现有方法主要分为基于输出信息的黑盒方法和基于内部表征的白盒方法。现有方法虽已取得一定进展，但仍存在两方面不足：其一，在通用幻觉检测领域中，对内部表征信息利用不充分，常局限于单层、单token或静态表示；其二，在RAG幻觉检测领域中，对外部检索知识信号利用不足，难以有效刻画模型内部与外部检索知识之间的冲突关系。

针对通用幻觉检测任务，本文提出ATScore，一种基于多方向表征建模的幻觉检测方法。该方法从大语言模型生成过程中的隐藏状态出发，构建统一的三维激活张量表示，并分别沿层方向与token两个方向进行链式表征建模，通过提取链长度、转折角度等几何统计特征，构造无需额外训练的白盒幻觉检测信号。相较于现有多数仅沿单一方向建模、仅利用局部表征的方法，ATScore能够在统一视角下更充分地利用跨层与跨token信息，挖掘内部表征演化过程中与幻觉相关的异常模式。实验结果表明，ATScore在多个公开数据集和多种评价指标上均取得了良好的检测效果，整体优于多种黑盒与白盒基线方法，表现出较好的稳定性、跨数据集泛化能力和跨场景适用性。

针对RAG幻觉检测任务，本文进一步提出SKID，一种基于检索知识冲突消解的幻觉检测方法。该方法首先利用稀疏自编码器将隐藏状态映射到高维稀疏特征空间，以获得更具语义可解释性的稀疏表征；随后通过互信息筛选出与幻觉判别最相关的特征维度，并结合检索文档对输出token进行过滤，重点关注未被外部知识直接支持的生成片段，以更直接地刻画并消解模型内部知识与外部检索知识之间的潜在冲突；最后采用广义加性模型实现对RAG幻觉的有效检测。实验结果表明，SKID在多个RAG数据集和多种评价指标上均具有较强竞争力，整体优于多类基线方法。进一步分析还表明，SKID不仅能够提升RAG幻觉检测性能，还能够通过稀疏特征的语义归纳解释模型被判定为幻觉的原因，从而兼顾检测准确性与结果可解释性。

\end{abstract}

%% ----------------------------------------------------------------------------
%%                              English Abstract
%% ----------------------------------------------------------------------------
\begin{englishabstract}{Large Language Models, Hallucination Detection, Internal Representations, Retrieval-Augmented Generation, Sparse Autoencoder}

In recent years, large language models have demonstrated strong capabilities in tasks such as question answering, text generation, and retrieval, yet they still tend to produce hallucinated content that is factually incorrect, logically flawed, or unsupported by evidence, thereby undermining the reliability and safety of such systems. The Retrieval-Augmented Generation (RAG) framework enhances model generation by introducing external retrieved knowledge sources and alleviates hallucination to some extent; however, it does not eliminate hallucinations entirely, and models may still generate content inconsistent with retrieved knowledge. The goal of hallucination detection is to effectively identify hallucinated content generated by large language models. Existing methods can be broadly categorized into black-box methods based on output information and white-box methods based on internal representations. Although existing studies have achieved some progress, they still suffer from two major limitations. First, in general hallucination detection, internal representation information is not fully exploited, and existing methods are often restricted to single-layer, single-token, or static representations. Second, in RAG hallucination detection, external retrieved knowledge signals are insufficiently utilized, making it difficult to effectively characterize the conflict between the model’s internal knowledge and external retrieved knowledge.

For the task of general hallucination detection, this thesis proposes ATScore, a hallucination detection method based on multi-directional representation modeling. Starting from the hidden states generated during the inference process of large language models, this method constructs a unified three-dimensional activation tensor representation, and further performs chain-based representation modeling along both the layer direction and the token direction. By extracting geometric statistical features such as chain length and turning angle, it constructs a white-box hallucination detection signal without requiring additional training. Compared with existing methods that mostly model representations along only a single direction and rely only on local representations, ATScore can make more sufficient use of cross-layer and cross-token information under a unified perspective, thereby capturing anomalous patterns related to hallucinations during the evolution of internal representations. Experimental results show that ATScore achieves strong detection performance on multiple public datasets and evaluation metrics, overall outperforming a variety of black-box and white-box baselines, and demonstrating good stability, cross-dataset generalization ability, and cross-scenario applicability.

For the task of RAG hallucination detection, this thesis further proposes SKID, a hallucination detection method based on retrieved knowledge conflict resolution. This method first maps hidden states into a high-dimensional sparse feature space by using a sparse autoencoder, so as to obtain sparse representations with better semantic interpretability. It then selects the feature dimensions most relevant to hallucination discrimination through mutual information, and filters output tokens in conjunction with retrieved documents, focusing on generated fragments that are not directly supported by external knowledge, so as to more directly characterize and resolve the potential conflict between the model’s internal knowledge and external retrieved knowledge. Finally, a generalized additive model is adopted to achieve effective RAG hallucination detection. Experimental results show that SKID is highly competitive on multiple RAG datasets and evaluation metrics, and overall outperforms various baseline methods. Further analysis also shows that SKID can not only improve the performance of RAG hallucination detection, but also explain why the model is judged as hallucinating through semantic induction over sparse features, thereby balancing detection accuracy and result interpretability.

\end{englishabstract}
